\section{Main}\label{sec:introduction}

% TODO: Figure 1 should illustrate: (a) disease-associated vs disease-relevant distinction,
% (b) the discovery cycle concept showing saturation patterns, (c) PIGEAN framework overview

% P1: The fundamental distinction - associated vs relevant
Drug targets with human genetic support are substantially more likely to succeed in clinical trials \cite{plenge_validating_2013,minikel_refining_2024}, yet fewer than 5\% of active drug programs leverage germline genetic data \cite{minikel_refining_2024}. This gap persists in part because the field conflates two fundamentally distinct concepts: \textit{disease-associated genes}, which carry genetic variants that reach statistical significance in population studies, and \textit{disease-relevant genes}, which causally influence disease when perturbed but may lack detectable genetic associations. Disease-associated genes are a subset of disease-relevant genes (Figure~\ref{fig:main}a), constrained by statistical power \cite{wang_statistical_2019}, the presence and frequency of functional variants \cite{gazal_linkage_2017}, and the difficulty of mapping noncoding signals to effector genes \cite{schipper_demystifying_2022,costanzo_realizing_2025}. The full set of disease-relevant genes---including both core mechanistic genes and peripheral network modifiers---is what matters for therapeutic development \cite{finan_druggable_2017,jhamb_pathway_2019}, yet we lack a principled framework for distinguishing relevance from association.

% P2: The discovery cycle problem - where are we for each trait?
This distinction has practical consequences for how we allocate research resources. For any given trait, GWAS discovery follows a characteristic cycle: early stages are limited by statistical power, but eventually the rate of new discoveries saturates even as sample sizes grow \cite{abdellaoui_15_2023}. At saturation, additional GWAS investment yields diminishing returns---further progress requires identifying disease-relevant genes through other means. Yet we currently have no systematic way to diagnose where a trait sits in this discovery cycle, leading to continued investment in ever-larger biobanks for traits that may have already saturated, while neglecting the development of better functional annotations that could unlock the remaining disease-relevant genes.

% P3: Non-genetic approaches and their limitations
Non-human-genetic approaches attempt to identify disease-relevant genes beyond those captured by GWAS. Experimental studies in animal \cite{seok_genomic_2013,widjaja_inhibiting_2019} and cellular \cite{ben-david_genetic_2018,canon_clinical_2019} systems establish causal gene-phenotype links but with uncertain human relevance. Observational studies using gene expression \cite{wu_integrating_2022,rodriguez_machine_2021} or epidemiological associations \cite{barter_cholCesteryl_2011,bhatt_cardiovascular_2019} provide human context but lack causal directionality. These sources populate biomedical databases widely used for drug discovery \cite{swinney_how_2011,dornbos_evaluating_2022} and can be efficiently queried with large language models \cite{toufiq_harnessing_2023,hu_evaluation_2025}. However, gene prioritization from these sources is prone to well-documented biases toward well-studied genes, incompleteness, and circularity \cite{gillis_assessing_2013,haynes_gene_2018,stoeger_large_2018,tan_caution_2025}.

% P4: Existing integration methods are heuristic
Ideally, human genetic associations would be combined with non-genetic evidence to identify disease-relevant genes more comprehensively. Expert synthesis achieves this for individual genes \cite{mountjoy_open_2022,karjalainen_genome-wide_2024}, but genome- and phenome-wide mapping requires scalable statistical methods. Such methods exist for GWAS locus-to-gene mapping, where functional annotations enriched for associations prioritize nearby genes \cite{pers_biological_2015,weeks_leveraging_2023}, and for rare disease diagnostics, where literature mining \cite{birgmeier_amelie_2020} and interaction networks \cite{smedley_next-generation_2015} prioritize patient variants. These approaches, however, rely on heuristics such as harmonic averaging \cite{buniello_open_2025} or regression on noisy training data \cite{weeks_leveraging_2023,schipper_prioritizing_2025}, lack explicit biological models, and do not distinguish association from relevance.

% P5: Our contribution - PIGEAN framework and discovery cycle diagnosis
Here, we introduce a Bayesian framework that separately quantifies \textit{direct genetic support} (the probability that a gene is disease-associated based on genetic data) and \textit{indirect genetic support} (the probability that a gene is disease-relevant based on shared functional annotations with associated genes). Indirect support is calibrated to the same scale as direct support, enabling principled combination while preserving the distinction between association and relevance. Applied across {{NUM_TOTAL_TRAITS_MOUSE_MSIGDB:,}} traits ({{NUM_COMMON_TRAITS_MOUSE_MSIGDB:,}} common, {{NUM_RARE_TRAITS_MOUSE_MSIGDB:,}} rare) and {{NUM_TOTAL_GENE_ANNOTATIONS:,}} gene annotations, we use this framework to diagnose discovery status across the phenome, finding that {{PERCENT_TRAITS_SATURATED}}\% of common diseases show evidence of GWAS saturation---where indirect support continues to identify disease-relevant genes while direct support has plateaued. We validate indirect support on three independent axes: it predicts held-out genetic associations, prioritizes drug targets with higher clinical success rates than existing resources, and enriches for causal genes in rare disease patients. This framework provides both a diagnostic for research resource allocation and a practical tool for identifying disease-relevant genes beyond the reach of GWAS alone.

% ANALYSIS IDEA: Consider adding a "headline number" for how many additional drug targets
% indirect support surfaces compared to direct-only approaches at equivalent success rates.
% This would strengthen the translational impact claim.

% TODO: PERCENT_TRAITS_SATURATED - Calculate from the ~1000 trait saturation analysis.
% Operational definition needed: e.g., traits where slope of indirect vs direct curve < 0.1
% at current sample size, or where indirect/direct ratio exceeds threshold.
